\documentclass[codesnippet]{jss}

\usepackage[T1]{fontenc}
\usepackage{orcidlink,lmodern}
\usepackage{booktabs}
\usepackage{amsmath}
\usepackage{amssymb}
\usepackage{graphicx}
\usepackage{url}
\usepackage{microtype}
\emergencystretch=3em
\raggedbottom
\DeclareFontShape{T1}{lmr}{b}{it}{<-> ssub * lmr/bx/it}{}

\newcommand{\class}[1]{`\code{#1}'}
\newcommand{\fct}[1]{\code{#1()}}
\urlstyle{tt}
\newcommand{\file}[1]{\path{#1}}

%% -- Article meta-information ------------------------------------------------

\author{Jason Turner\\Independent Researcher}
\Plainauthor{Jason Turner}

\title{\pkg{pycorpdiff}: Comparative Corpus Analysis for Modern \proglang{Python} Workflows}
\Plaintitle{pycorpdiff: Comparative Corpus Analysis for Modern Python Workflows}
\Shorttitle{\pkg{pycorpdiff}: Comparative Corpus Analysis in \proglang{Python}}

\Abstract{
Corpus linguistics, digital humanities, and computational social
science routinely require *comparing* text corpora --- across
political alignments, across time, across discourse contexts. The
existing tooling is fragmented across the corpus-linguistics
ecosystem (\proglang{R}'s \pkg{quanteda}, SketchEngine),
visualisation-first packages (\pkg{Scattertext}, \pkg{BERTopic}), and
the \proglang{Python} NLP stack (\pkg{nltk}, \pkg{spaCy},
\pkg{gensim}, \pkg{sentence-transformers}). No \proglang{Python}
package has unified the comparative workflow under a single coherent
API.

This paper presents \pkg{pycorpdiff}, a \proglang{Python} package
that provides a comparative layer over the existing PyData and NLP
stacks. The package exposes three public verbs ---
\fct{compare(a, b)}, \fct{track(c, term)}, and
\fct{compare.before\_after()} --- that consolidate keyness analysis
(Dunning log-likelihood with LogRatio effect sizes, dispersion
sanity checks, multiple-comparison correction), collocation shift
(window-based co-occurrence with logDice, PMI, $t$-score, MI\textsuperscript{3}),
temporal trajectories with Wilson score confidence intervals,
changepoint detection, interrupted time series, and embedding-based
semantic shift via averaged contextual representations. Every
analytical result carries its evidence: each Result object exposes
\fct{.explain(term)} which returns KWIC concordance lines from both
source corpora.

\pkg{pycorpdiff} is positioned as a comparative layer, not a general
NLP framework. It does not reinvent tokenisation, embeddings, or
topic modelling --- those concerns plug in via two \class{typing.Protocol}
extension points (\code{Tokenizer} and \code{Embedder}). The base
install depends only on \pkg{numpy}, \pkg{pandas}, \pkg{scipy}, and
\pkg{pyarrow}; visualisation, semantic, temporal-statistics, and
multilingual functionality are opt-in extras.

We describe the package's three-layer architecture, the statistical
defaults that drive each analytical method, and a worked example on
a synthetic two-frame discourse corpus. A replication archive
accompanies the manuscript so every figure and table in the paper
can be reproduced from the published code.
}

\Keywords{comparative corpus analysis, keyness, collocations,
semantic change, diachronic NLP, digital humanities,
computational social science, reproducible research, \proglang{Python}}

\Plainkeywords{comparative corpus analysis, keyness, collocations,
semantic change, diachronic NLP, digital humanities,
computational social science, reproducible research, Python}

\Address{
  Jason Turner\\
  Independent Researcher\\
  E-mail: \email{jason.s.turner@gmail.com}\\
  URL: \url{https://github.com/jasonsturner/pycorpdiff}
}

\begin{document}

\section{Introduction}\label{sec:intro}

[TODO: motivate the comparative-corpus problem. Researchers in
corpus linguistics, digital humanities, computational social
science, communication studies, and discourse analysis frequently
need to answer: how does corpus A differ from corpus B?\ how has
language around \emph{X} evolved over time?\ what discourse
changed before vs.\ after a known event?\ Existing tooling is
fragmented across at least four ecosystems: classical corpus
linguistics (\proglang{R}'s \pkg{quanteda}, the Sketch Engine
platform), visualisation-first systems (\pkg{Scattertext}), modern
\proglang{Python} NLP (\pkg{nltk}, \pkg{spaCy}, \pkg{gensim},
\pkg{sentence-transformers}), and temporal/time-series tooling
(\pkg{statsmodels}, \pkg{ruptures}, \pkg{sktime}). Each does its
job well, none provides the comparative-analysis layer.]

[TODO: state the contribution. \pkg{pycorpdiff} is the missing
comparative layer. Three verbs (\fct{compare}, \fct{track},
\fct{compare.before\_after}) consolidate keyness, collocations,
semantic shift, temporal trajectories, changepoint detection, ITS,
and explainable evidence into a single notebook-native API. The
package is positioned as orchestration: it integrates the existing
ecosystem rather than reimplementing it.]

\section{Related work}\label{sec:related}

\subsection{Corpus linguistics and keyness}

[TODO: Scott (1997) on WordSmith; Kilgarriff on corpus comparison;
Dunning (1993); Gabrielatos \& Marchi on keyness methodology;
Hardie (2014) on LogRatio; Wilson (2013) on Bayes-factor keyness;
Rayson's LL Wizard; Juilland (1964) and Gries (2008) on dispersion.]

\subsection{Sketch Engine and the comparative tradition}

[TODO: Kilgarriff et al.\ on SketchEngine; word sketches, sketch
diff, multilingual comparison. Note: \pkg{pycorpdiff} is not a
SketchEngine clone --- it deliberately stays in \proglang{Python}'s
data-frame idiom and integrates with the \proglang{Python} embedding
ecosystem rather than reproducing SketchEngine's hosted-platform model.]

\subsection{Semantic change}

[TODO: Hamilton et al.\ (2016) diachronic word embeddings;
Procrustes alignment for independent per-corpus spaces;
Giulianelli et al.\ (2020) on averaged contextual embeddings
under shared encoders like BERT/SBERT.]

\subsection{R ecosystem}

[TODO: \pkg{quanteda} (Benoit et al.\ 2018) and \pkg{tidytext} (Silge
\& Robinson 2016). Note shared idioms (tidy data frames, pipelines)
and differences (\pkg{quanteda}'s document-feature matrix vs.\
\pkg{pycorpdiff}'s dataframe-of-documents).]

\subsection{Python ecosystem}

[TODO: \pkg{nltk}, \pkg{spaCy}, \pkg{gensim}, \pkg{sentence-transformers},
\pkg{Scattertext}, \pkg{BERTopic}. Each is best-of-breed in its
specific niche; none is a comparative-analysis layer.]

\section{Design}\label{sec:design}

\pkg{pycorpdiff} is structured as three concentric layers, each
with a single responsibility:

\begin{itemize}
  \item \textbf{Layer 1 (data plumbing).} \class{Corpus},
    \class{CorpusSlice}, \class{Tokenizer} (a \class{typing.Protocol}),
    \class{Embedder} (a \class{typing.Protocol}), and the I/O
    readers for txt, CSV, parquet, and \pkg{pandas}.
  \item \textbf{Layer 2 (methods).} Pure mathematical functions
    that take dataframes and return dataframes.
    \fct{log\_likelihood}, \fct{log\_ratio},
    \fct{percent\_diff}, \fct{bayes\_factor},
    \fct{juilland\_d}, \fct{dispersion\_dp},
    \fct{benjamini\_hochberg}; \fct{logdice}, \fct{pmi},
    \fct{t\_score}, \fct{mi\_three}, \fct{collocation\_shift};
    \fct{semantic\_shift}, \fct{neighborhood\_drift},
    \fct{procrustes\_align}; \fct{detect\_changepoints},
    \fct{interrupted\_time\_series}.
  \item \textbf{Layer 3 (public verbs and results).}
    \fct{compare}, \fct{track}, \fct{compare.before\_after},
    plus the four \class{Result} dataclasses
    (\class{KeynessResult}, \class{CollocationShiftResult},
    \class{SemanticShiftResult}, \class{TemporalTrajectory}).
\end{itemize}

[TODO: expand on the architectural invariants. Layer 2 never
imports Layer 3. Layer 1 is backend-pluggable. Tokenizer and
Embedder are \class{typing.Protocol}, not classes --- spaCy /
Stanza / jieba / fugashi adapters are one-line wrappers. Every
Result implements the same informal contract: \fct{.to\_df()},
\fct{.plot()}, \fct{.explain()}, \fct{.summary()}.]

\section{API}\label{sec:api}

\subsection{Constructing a corpus}

[TODO: code listing of \fct{from\_dataframe}, \fct{read\_csv},
\fct{read\_parquet}, slicing, \fct{by\_time}. Reference
Section~\ref{sec:design} for the type structure.]

\subsection{Keyness}

[TODO: \fct{compare(a, b).keyness} with effect-size and
dispersion options. Show the volcano plot and the top-N bar.
Reference Hardie's LogRatio and Wilson's BF.]

\subsection{Collocation shift}

[TODO: \fct{compare(a, b).collocation\_shift(target)}; logDice
default, alternative measures, Laplace smoothing for
absent-on-one-side collocates.]

\subsection{Temporal trajectories}

[TODO: \fct{track(corpus, term).over\_time} with Wilson CIs.
Show the line plot with the CI band.
\fct{TemporalTrajectory.changepoints} and \fct{.interrupted\_time\_series}
for follow-up analyses.]

\subsection{Semantic shift}

[TODO: \fct{compare(a, b).semantic\_shift(target, embedder)} via
averaged contextual embeddings. SBERT as the default;
\class{HashEmbedder} for deterministic testing; Procrustes
alignment for Hamilton-style independent spaces.]

\subsection{Explainability}

[TODO: \fct{KeynessResult.explain} and
\fct{CollocationShiftResult.explain} return \class{ConcordanceResult}
with KWIC lines from both source corpora. Every ranked term
carries its evidence by construction.]

\section{Worked example}\label{sec:example}

[TODO: a single end-to-end run on a synthetic two-frame migrant-
discourse corpus. The replication script in
\file{paper/replication/reproduce.py} regenerates every figure
and table in this section. Walk through: keyness, collocation
shift, temporal trajectory with engineered 2016 event,
changepoint detection, ITS, semantic shift via \class{HashEmbedder}
(for byte-reproducibility in CI).]

\section{Statistical defaults}\label{sec:defaults}

[TODO: one paragraph per metric explaining what the default
does and why. Keyness: signed G\textsuperscript{2} primary,
LogRatio as effect size, BH adjustment. Collocations: logDice
default with Laplace smoothing. Trajectories: Wilson score CIs
at 95\%. Changepoints: PELT with rbf cost, BIC-style penalty.
ITS: segmented regression via \pkg{statsmodels.OLS}. Semantic
shift: averaged contextual embeddings, no alignment by default
under shared-model encoders. All defaults are configurable.]

\section{Reproducibility and testing}\label{sec:reproducibility}

[TODO: the package ships 193 tests including known-answer tests
against Rayson's LL Wizard (G\textsuperscript{2} = 182.07 for
12000/1M vs 10000/1M), Newcombe's Wilson CI tables, Schoenemann's
Procrustes invariants, and Hardie-recommended LogRatio behaviour.
Hypothesis-based property tests guard swap-symmetry,
non-negativity of $|G^2|$, and proportional-corpora-equal-rate
invariants. The tutorial notebook is executed in CI on every push
so it cannot silently drift. Coverage exceeds 90\% on the keyness
and collocation core.]

\section{Conclusion}\label{sec:conclusion}

[TODO: \pkg{pycorpdiff} provides the missing comparative layer
over the \proglang{Python} corpus-analysis ecosystem. Future work
includes a polars backend, the multilingual demo corpora,
real-world worked examples on UK Hansard and CORD-19, and the
plug-in path for trained domain embeddings.]

\section*{Computational details}

The results in this paper were obtained using \proglang{Python}
3.12 with \pkg{pycorpdiff}~0.1.0a0, \pkg{numpy}~$\ge$~1.24,
\pkg{pandas}~$\ge$~2.0, \pkg{scipy}~$\ge$~1.11,
\pkg{pyarrow}~$\ge$~14, \pkg{altair}~$\ge$~5,
\pkg{sentence-transformers}~$\ge$~2.2, \pkg{ruptures}~$\ge$~1.1,
and \pkg{statsmodels}~$\ge$~0.14. \pkg{pycorpdiff} is freely
available from the project repository at
\url{https://github.com/jasonsturner/pycorpdiff} under the MIT
license.

\section*{Acknowledgments}

[TODO]

% While the manuscript is still a skeleton (\TODO placeholders in every
% section), no \cite{} calls exist yet — bibtex would warn "no citations".
% \nocite{*} forces every entry in references.bib to appear in the
% rendered bibliography for now; remove it once the prose pass cites
% references explicitly via \cite{}.
\nocite{*}
\bibliography{references}

\end{document}
